\documentclass[10pt,a4paper]{amsart}
\usepackage[margin=19mm]{geometry}
\usepackage{amsmath,amssymb,mathtools}
\usepackage[scr=rsfs,cal=euler]{mathalfa}
\usepackage{xfrac}
\usepackage[unicode,hidelinks,bookmarks=false]{hyperref}
\newcommand{\ie}{i.e.\ }
\newcommand{\cf}{cf.\ }
\newcommand{\resp}{respectively\ }
\newcommand{\Subsection}{\subsection}
\providecommand{\nobreakdash}{\nobreak\hbox{-}}
\setlength{\emergencystretch}{2em}
\multlinegap0pt
\usepackage{bm}
\usepackage{bbm}
\usepackage{dsfont}
\usepackage{xspace}
\usepackage{cancel}
\usepackage{mathtools}
\usepackage{amsthm}
\makeatletter
\@ifundefined{theorem}{\newtheorem{theorem}{Theorem}}{}
\makeatother
\newcommand{\on}{\operatorname}
\title[Duality and linearization for p-adic lie groups]{Duality and linearization for p-adic lie groups}
\author{Dustin Clausen}
\date{Source: arXiv:2506.18174v1 (CC BY 4.0)}
\begin{document}
\maketitle

\noindent\emph{Source and licence.} Duality and linearization for p-adic lie groups. Version arXiv:2506.18174v1 (CC BY 4.0), distributed under the Creative Commons Attribution 4.0 International licence, as recorded in the arXiv licence metadata for this exact version and shown on the abstract page https://arxiv.org/abs/2506.18174v1. Full rights notice: licenses/arxiv-2506.18174-CC-BY-4.0.txt.\par\smallskip

\noindent\textit{Editorial scope.} Counted unit: the Theorem whose source label is \texttt{heightn} in the article (source0.tex lines 2877--2884, source label \texttt{heightn}), delivered together with the article's own complete original proof of it (source0.tex lines 2885--2915, one \texttt{proof} environment of 31 lines). No mathematics is written, repaired or re-derived: statement and proof are the article's own text line for line. Required context retained uncounted: none. Every definition, display and label of those lines is retained unchanged. Omitted from this package and not redistributed: 1--2876, 2916--3569. Nothing inside a retained excerpt is omitted or altered: every retained body is delivered exactly as the article's lines are written, so no omission receipt of any kind accompanies this package. Citations: the retained excerpts carry no citation command at all, so no bibliography entry of the article is transcribed and no reference is redistributed. Reference adaptation: none was needed and none was made; every reference inside the delivered unit resolves inside it, so no unresolved reference or citation is produced. Printed slips kept literal: no printed word, symbol, hypothesis or number of the counted statement or of its proof was altered. Layout and class: the article's own class is \texttt{article} with packages including fullpage, amsmath, amsfonts, amsthm, amssymb, epsfig, biblatex, xy, enumitem, xcolor, hyperref, graphicx, tikz-cd, float; the wrapper is the lane's pinned \texttt{amsart} editorial wrapper at 10pt on A4, which restates the theorem-like environments the retained text needs and adds the article's own macro definitions that the retained text needs (\texttt{\textbackslash on}), with extra wrapper packages (bm, bbm, dsfont, xspace, cancel, mathtools). The delivered numbering is the unit's own. Assets: no retained span contains a figure environment, a graphics-inclusion command, a PDF-page inclusion, a table or a TikZ picture; no image file is copied into this package, the package declares no asset dependency and no image file is redistributed with it. Licence: version arXiv:2506.18174v1 of this article is distributed under the Creative Commons Attribution 4.0 International licence, as recorded in the arXiv licence metadata for this exact version and shown on the abstract page https://arxiv.org/abs/2506.18174v1. A full rights notice is delivered at licenses/arxiv-2506.18174-CC-BY-4.0.txt. Characters: every retained excerpt is pure ASCII, so the delivered engine, XeLaTeX with its default Latin Modern text font, reports no missing character.
\begin{theorem}\label{heightn}
Let $n\geq 1$.  Then:
\begin{enumerate}
\item Every one-dimensional formal group $\widehat{G}$ of exact height $n$ over an $\mathbb{F}_p$-algebra $R$ is pro-etale locally isomorphic to $\widehat{H}_n$.
\item The pro-etale sheaf of rings $\on{End}(\widehat{H}_n)$ of endomorphisms of $\widehat{H}_n$ over $\mathbb{F}_p$, when base-changed to $\mathbb{F}_{p^n}$, is represented by the pro-constant ring scheme on the profinite ring $\mathcal{O}_n$.
\item The identification in 2 is equivariant with respect to $\on{Gal}(\mathbb{F}_{p^n}/\mathbb{F}_p)$, which acts on $\mathcal{O}_n$ since this ring was functorially constructed in terms of $\mathbb{F}_{p^n}$.
\end{enumerate}
\end{theorem}

\begin{proof}
First consider 1.  Working Zariski locally on $\on{Spec}(R)$ we can assume that the Cartier module of $\widehat{G}$ is presented by a single generator $e$ with the single relation
$$Fe = \sum_{d\geq 0} V^{d+n-1} [c_d]e$$
where $c_d\in R$ for all $d\geq 0$ and $c_0\in R^\times$.  We need to find, pro-etale locally, a new $V$-basis vector $e'$ such that $Fe'=V^{n-1}e'$.  Write
$$e' = \sum_{k\geq 0} V^k[b_k]e.$$
The condition that $e'$ be a $V$-basis is $b_0\in R^\times$.  Furthermore, using relations in the Cartier ring and the above expression for $Fe$ we compute that
$$Fe' = \sum_{d,k\geq 0} V^{d+k+n-1}[\varphi^{d+n}(b_k)\cdot c_d]e$$
whereas
$$V^{n-1}e' = \sum_{k\geq 0} V^{k+n-1}[b_k]e.$$
Equating coefficients, we find that $e'$ does what we want if and only if
$$b_0 = \varphi^n(b_0)\cdot c_0,$$
$$b_1 = \varphi^{n+1}(b_0)\cdot c_1 + \varphi^n(b_1)\cdot c_0,$$
$$b_2 = \varphi^{n+2}(b_0)\cdot c_2 + \varphi^{n+1}(b_1)\cdot c_1 + \varphi^n(b_2)\cdot c_0,$$
etc.  Divide through by $c_0$.  Then when we try to inductively solve for the $b_k$, we find that each equation amounts to a monic polynomial that $b_k$ should be a root of, and that this monic polynomial has derivative the constant polynomial on the unit $-c_0^{-1}$, hence describes a finite etale extension.  To enforce that $b_0$ be a unit, we can divide the first equation by $b_0$ and see that we're still left with a finite etale equation $b_0^{p^n-1}=c_0^{-1}$ for $b_0$, but one which forces $b_0$ to be a unit.  Thus at the top of this tower of finite etale extensions we get an isomorphism $\widehat{G}\simeq \widehat{H_n}$, as required.

More precisely, what the above argument showed was that the functor on $R$-algebras $\on{Hom}(\widehat{H_n},\widehat{G})$ parametrizing the set of homomorphisms $\widehat{H_n}\to\widehat{G}$ is canonically the inverse limit
$$\on{Hom}(\widehat{H_n},\widehat{G})\simeq \varprojlim_i \on{Hom}_i(\widehat{H_n},\widehat{G})$$
where $\on{Hom}_i(\widehat{H_n},\widehat{G})$ parametrizes homomorphisms of the associated Cartier modules mod $V^i$ (so in particular $\on{Hom}_0(\widehat{H_n},\widehat{G})=\on{Spec}(R)$), and
$$\on{Hom}_{i+1}(\widehat{H_n},\widehat{G})\to \on{Hom}_i(\widehat{H_n},\widehat{G})$$
is finite etale and surjective for all $i\geq 0$, described by the sequence of separable monic polynomials discussed above; moreover the subfunctor $\on{Isom}(\widehat{H_n},\widehat{G})\subset \on{End}(\widehat{H_n},\widehat{G})$ parametrizing isomorphisms comes by pullback from $\on{Isom}_1(\widehat{H_n},\widehat{G})$, which is the finite etale extension of $\on{Spec}(R)$ described by $c_0X^{p^n-1}=1$ inside $\on{Hom}_1(\widehat{G},\widehat{H_n})=\on{Spec}(R[X]/(c_0X^{p^n}-X))$.

Now, moving on to 2, we specialize the above to $\widehat{G}=\widehat{H}_n$. Then $c_0=1$ and $c_d=0$ for $d>0$, so the above equations just amount to $b_k=\varphi^n(b_k)$ for all $k\geq 0$.  The polynomial $x^{p^n}-x$ splits completely over $\mathbb{F}_{p^n}$, so all the finite etale extensions appearing above are constant over $\mathbb{F}_{p^n}$, proving the pro-constancy of $\on{End}(\widehat{H}_n)$ over $\mathbb{F}_{p^n}$.  To finish the proof of 2, it suffices to identify each ring $\on{End}(\widehat{H}_n)_i(\mathbb{F}_{p^n})$ with $\mathcal{O}_n/\pi^i$, compatibly in $i$, and to prove 3 we have to make this identification functorial in the choice of $\mathbb{F}_{p^n}$.

By the above analysis, if $M$ denotes the Cartier module of $\widehat{H}_n$, then the set of endomorphisms of $M/V^i M$ over $\mathbb{F}_{p^n}$ is in bijection with the set of sequences
$$b_0,b_1,\ldots,b_{i-1}$$
of elements of $\mathbb{F}_{p^n}$, where to such a sequence we assign the endomorphism $f$ characterized uniquely by
$$f(e) = ([b_0] + V[b_1] + \ldots + V^{i-1}[b_{i-1}])e.$$
This is also in bijection with elements of the Cartier ring (mod $V^i$) of the form
$$[b_0] + V[b_1] + \ldots + V^{i-1}[b_{i-1}].$$
Using the injective homomorphism of $W(\mathbb{F}_{p^n})$ into the Cartier ring, assigning to a Witt vector with Witt coordinates $(a_0,a_1,\ldots)$ the element $\sum_{d\geq 0} V^d [a_d] F^d$ of the Cartier ring, we easily calculate that the ring structure on endomorphisms, via this parametrization, indeed corresponds to the ring structure on $\mathcal{O}_n/\pi^i$, with $\pi$ corresponding to $V$ (which, as an endomorphism of $\widehat{H}_n$, corresponds to the relative Frobenius $t\mapsto t^p$).  The key relation $V^n=p$ amounts to the definition of $\widehat{H}_n$.
\end{proof}

\end{document}
